\section{课程收尾：把这些知识拼成一个系统}

\begin{figure}[H]
\centering
\begin{tikzpicture}[node distance=0.95cm, every node/.style={font=\small, align=center}]
\node[draw, rounded corners, fill=black!10, minimum width=4.6cm, minimum height=0.8cm] (a) {build};
\node[draw, rounded corners, fill=black!18, below=of a, minimum width=4.6cm, minimum height=0.8cm] (b) {measure};
\node[draw, rounded corners, fill=black!26, below=of b, minimum width=4.6cm, minimum height=0.8cm] (c) {scale};
\node[draw, rounded corners, fill=black!34, below=of c, minimum width=4.6cm, minimum height=0.8cm] (d) {align};
\draw[-{Latex[length=3mm]}, thick] (a) -- (b);
\draw[-{Latex[length=3mm]}, thick] (b) -- (c);
\draw[-{Latex[length=3mm]}, thick] (c) -- (d);
\end{tikzpicture}
\caption{课程给出的最终闭环：build, measure, scale, align}
\end{figure}

\begin{knowledgebox}{最后应该带走的三句话}
\begin{enumerate}
    \item 这门课不是教你“调用模型”，而是教你“构建模型”。
    \item 真正有价值的是能随着规模继续成立的算法和系统。
    \item 理解语言模型，必须把数据、系统和模型一起看。
\end{enumerate}
\end{knowledgebox}

\begin{warningbox}{最后一句话}
不要只做会调用模型的人，要做知道模型为什么能工作的那个人。
\end{warningbox}

